\chapter{Introduction}
\label{sec:intro}

It is a widely accepted truth in the world of music production and recording that spatial information in audio is crucial to  listeners' sense of immersion and connection with the music. To capture that spaciousness, stereophonic recording and reproduction has been the standard format since it gained popularity in the 1950s \cite{zagorski-thomasBloomsburyHandbookMusic2020,lanierHIFITWOCHANNELCOMMOTION1957}. The spatial analysis of subsequent decades worth of stereo recordings presents a unique opportunity to study human musicality, as each recording reflects and encodes the creative choices made by the musicians and recordists involved.

Within the music community it is common to speak about spaciousness in a qualitative sense, using terms such as \enquote{stereo image} and \enquote{stereo field} \footnote{The terms \q{stereo image} and \q{stereo field} are freely interchangeable and, as is normal in the music community, I use them both throughout this thesis.} to encompass all audible spatial aspects. In academia however -- and specifically in the \gls{mir} community -- there have been several explicit attempts to translate these discussions into a quantitive domain \cite{tzanetakisStereoPanningFeatures2007,sarroffMeasurementsSpaciousnessStereophonic2008,tardieuProductionEffectAudio2011}. All were performed in an era when musical datasets were \mbox{hand-curated/-labelled}, fairly small (\numrange{1000}{2000} tracks), and, most notably, private. In recent years new musical corpora have become available, e.g. MTG-Jamendo \cite{bogdanovMTGJamendoDatasetAutomatic2019} and \gls{fma} \cite{defferrardFMADatasetMusic2017}, which are both orders of magnitude larger than their predecessors, containing approximately \num{55000} and \num{100000} tracks respectively. In addition, computing techniques in \gls{ml} and \gls{dl} have progressed to the point where it is now possible to learn extremely high-level features on these vast datasets, thereby making sense of the human creativity represented in the recordings. 

A typical use case for spatial descriptors that has emerged both in academia and industry is classification, and in particular, auto-tagging of musical genre/performer \cite{tzanetakisStereoPanningInformation2010,ziemerGoniometersArePowerful2023}. Such features can be used in music similarity and recommendation tasks \cite{gomezMusicSimilaritySystems2008,ziemerNovelRecordingStudio2020}. It has also been shown that stereo field classification of subsections within a larger work can act as a precursor to other tasks such as segmentation, audio enhancement, upmixing, and source separation \cite{barryMusicStructureSegmentation2007,harmaStereoAudioClassification2011,harmaClassificationTimefrequencyRegions2011}. The latter is very relevant to the audio distribution industry, which is currently experimenting with newer audio formats such as ambisonic and object-based audio. While these new formats offer exciting opportunities for both distributors and listeners, there remains a huge back-catalogue of stereophonic recordings that need to remain playable on modern reproduction equipment. A detailed  understanding of their spatial content would therefore be beneficial.

Historically, binaural audio (which is a subclass of stereo) has been of particular interest because its relevance to human hearing. Thus, it is extremely widely studied, and the underlying fundamental techniques are well-established \cite{blauertTechnologyBinauralListening2013}. Yet there remains a gap in current research regarding direct analysis of raw music recordings, and in particular, describing the stereo image of said recordings at a high level. Through the work described in this thesis, I aim to fill that gap. 

I firstly demonstrate use cases where the spatial analysis of stereo audio has the opportunity to solve real-world problems in audio transmission, dissemination, and reproduction. Accordingly, I gather existing objective methods to describe stereo image and assess their abilities to detect the spatial aspects of a variety of musical audio sources.\footnote{Throughout this thesis I use the word \q{objective} quite frequently. The reader is directed to \cref{ch:methodology} for a full discussion of the philosophical implications, but in essence I use the term to infer that the signal processing techniques discussed are deterministic and applicable to any musical recording.} 